\documentclass[11pt]{article}
\usepackage[utf8]{inputenc}
\usepackage[T1]{fontenc}
\usepackage{amsmath,amssymb,amsthm}
\usepackage{geometry}
\geometry{margin=2.5cm}

\newtheorem{theorem}{Théorème}
\newtheorem{proposition}{Proposition}
\newtheorem{lemma}{Lemme}
\newtheorem{corollary}{Corollaire}
\newtheorem{definition}{Définition}
\newtheorem{hypothesis}{Hypothèse}
\theoremstyle{remark}
\newtheorem{remark}[theorem]{Remarque}

\newcommand{\Inc}{\ensuremath{\mathrm{Inc}}}
\newcommand{\Exc}{\ensuremath{\mathrm{Exc}}}
\newcommand{\vbar}{\bar{v}}

\begin{document}

\begin{center}
{\LARGE\bfseries Renforcement du Théorème 4 en équivalence}\\[4pt]
{\large Suite à la suggestion de Florent (analogue au Corollaire 2)}
\end{center}
\vspace{1em}

Florent a remarqué que le Corollaire 2 (non-contradiction) est en fait une
équivalence, pas seulement une implication : $(a\le c$ et $d\le b) \iff
(\Delta+\bar\Delta\le0$ pour toute loi$)$. Il propose la même chose pour le
Théorème 4 (identification). On montre ici que c'est vrai, en identifiant
précisément ce qui est trivial et ce qui demande une preuve nouvelle.

\section{Énoncé renforcé}

\begin{theorem}[Identification -- version équivalence]\label{thm:identif-equiv}
Sous l'Hypothèse~1, supposons $Y=Y^\star$ presque sûrement, la loi jointe de
$(X_1,\ldots,X_n)$ étant quelconque. Les trois assertions suivantes sont
équivalentes :
\begin{enumerate}
\item[(i)] Pour tout $j\in J^+$ : $aA_j+dD_j>cB_j$ ; pour tout $j\in J^-$ :
$aB_j+dC_j>cA_j$ ; pour tout $k\notin J^+\cup J^-$ : $aA_k+dD_k<bC_k+cB_k$
\emph{et} $aB_k+dC_k<cA_k+bD_k$.
\item[(ii)] $\Delta_j>0$ et $\bar\Delta_j<0$ pour $j\in J^+$ ; $\bar\Delta_j>0$
et $\Delta_j<0$ pour $j\in J^-$ ; $\Delta_k<0$ et $\bar\Delta_k<0$ pour
$k\notin J^+\cup J^-$.
\item[(iii)] $\displaystyle\lim_{N\to\infty}\limsup_{t\to\infty}P(C_t\ne C_{\mathrm{true}})=0$.
\end{enumerate}
\end{theorem}

\section{(i) $\iff$ (ii) : une identité, rien de nouveau}

\begin{proof}[Preuve de (i)$\iff$(ii)]
Pour $j\in J^+$ : $Y^\star=1\Rightarrow X_j=1$ (définition de $C_{\mathrm{true}}$),
donc $C_j=0$. Dans $\Delta_j=aA_j+dD_j-bC_j-cB_j$, le terme $bC_j$ s'annule :
$\Delta_j=aA_j+dD_j-cB_j$. Donc ``$\Delta_j>0$'' et ``$aA_j+dD_j>cB_j$'' sont
la \emph{même} affirmation, pas deux affirmations à relier. Sous
l'Hypothèse~1, le Corollaire~2 donne déjà $\Delta_j>0\Rightarrow\bar\Delta_j<0$ ;
la clause ``et $\bar\Delta_j<0$'' de (ii) est donc automatique, pas une
condition supplémentaire à vérifier. Symétriquement pour $j\in J^-$ (avec
$D_j=0$). Pour $k\notin J^+\cup J^-$, aucune simplification logique
n'intervient : $\Delta_k<0$ et $aA_k+dD_k<bC_k+cB_k$ sont, par définition
même de $\Delta_k$, la même inégalité ; idem pour $\bar\Delta_k$. Donc (i) et
(ii) sont mot pour mot la même chose, une fois les définitions substituées.
\end{proof}

\section{(ii) $\Rightarrow$ (iii) : déjà prouvé}

C'est exactement le contenu du Théorème 4 original combiné au Théorème 3
(concentration, version corrigée pour les cas absorbants) : si chaque
automate a la dérive de signe requis, $C^\star=C_{\mathrm{true}}$, et le
Théorème 3 donne la limite double annoncée. Rien à ajouter ici.

\section{(iii) $\Rightarrow$ (ii) : la partie nouvelle}

C'est ici la vraie contribution de la suggestion de Florent. On procède par
contraposée : on montre que si (ii) est en défaut pour \emph{un seul}
littéral, alors (iii) est faux -- la limite double ne vaut pas $0$, elle
vaut $1$ (ou au moins $1/2$ dans un cas limite).

\begin{lemma}\label{lem:witness}
S'il existe $j_0\in J^+$ tel que $\Delta_{j_0}\le0$, alors
$\lim_{N\to\infty}\limsup_{t\to\infty}P(C_t\ne C_{\mathrm{true}})=1$.
\end{lemma}

\begin{proof}
Deux cas, selon le signe de $\Delta_{j_0}$.

\emph{Cas $\Delta_{j_0}<0$.} Alors $\rho_{j_0}=p_{j_0}^+/p_{j_0}^-<1$. Par le
Théorème 1, la loi stationnaire de $v_{j_0}$ vérifie
$\pi(\Inc)=\rho_{j_0}^N/(1+\rho_{j_0}^N)$, donc
$\pi(\Exc)=1/(1+\rho_{j_0}^N)$. Comme $C_{\mathrm{true}}$ exige $x_{j_0}$
inclus (car $j_0\in J^+$), tout état $C_t$ où $v_{j_0}$ est en $\Exc$ vérifie
$C_t\ne C_{\mathrm{true}}$. Donc, par le théorème ergodique (Théorème 1) :
\[
  \limsup_{t\to\infty}P(C_t\ne C_{\mathrm{true}})
  \ge \limsup_{t\to\infty}P(v_{j_0,t}=\Exc) = \pi(\Exc) = \frac1{1+\rho_{j_0}^N}.
\]
Comme $\rho_{j_0}<1$, $\rho_{j_0}^N\to0$ quand $N\to\infty$, donc
$\frac1{1+\rho_{j_0}^N}\to1$. D'où
$\lim_{N\to\infty}\limsup_{t\to\infty}P(C_t\ne C_{\mathrm{true}})\ge1$, donc
$=1$ (une probabilité ne dépasse pas $1$).

\emph{Cas $\Delta_{j_0}=0$ (cas limite).} Sous l'Hypothèse de généricité
(pas $p^+=p^-=0$), on a alors $p_{j_0}^+=p_{j_0}^->0$, donc $\rho_{j_0}=1$.
Par le Théorème 1 (cas $\rho=1$), la loi stationnaire est uniforme :
$\pi(k)=1/(2N)$ pour tout $k$, donc $\pi(\Exc)=\pi(\Inc)=1/2$ --
\emph{exactement}, sans dépendre de $N$. Le même argument que ci-dessus donne
\[
  \limsup_{t\to\infty}P(C_t\ne C_{\mathrm{true}}) \ge \pi(\Exc) = \frac12,
  \qquad\text{et ce pour tout } N,
\]
donc $\lim_{N\to\infty}\limsup_{t\to\infty}P(C_t\ne C_{\mathrm{true}})\ge\frac12\ne0$.
\end{proof}

\begin{remark}
Le cas absorbant ($p_{j_0}^-=0$, $p_{j_0}^+>0$ mais dans le \emph{mauvais}
sens -- ce qui ne peut se produire ici que si en fait $\Delta_{j_0}>0$, donc
hors du champ de ce lemme) ne pose pas de problème symétrique : si
$\Delta_{j_0}\le0$ correspond à un cas absorbant (i.e. $p_{j_0}^+=0$),
l'automate se bloque en $\Exc$ en temps fini (Corollaire 1), ce qui donne
directement $\limsup_t P(C_t\ne C_{\mathrm{true}})=1$ pour tout $N$, encore
plus fort que la conclusion du Lemme~\ref{lem:witness}.
\end{remark}

\begin{lemma}\label{lem:witness-neg}
S'il existe $j_0\in J^-$ tel que $\bar\Delta_{j_0}\le0$, ou $k_0\notin
J^+\cup J^-$ tel que $\Delta_{k_0}\ge0$ ou $\bar\Delta_{k_0}\ge0$, alors de
même $\lim_{N\to\infty}\limsup_{t\to\infty}P(C_t\ne C_{\mathrm{true}})\in\{\tfrac12,1\}\ne0$.
\end{lemma}

\begin{proof}
Identique au Lemme~\ref{lem:witness}, en remplaçant $v_{j_0}$ par
$\vbar_{j_0}$ (pour $j_0\in J^-$, où $C_{\mathrm{true}}$ exige $\vbar_{j_0}$
en $\Inc$), ou par $v_{k_0}$ (resp. $\vbar_{k_0}$) pour $k_0$ non pertinent
(où $C_{\mathrm{true}}$ exige $\Exc$ des deux côtés). Le même calcul
(Théorème 1, cas $\rho<1$ absent ici puisqu'on suppose le signe manquant
dans l'autre sens -- $\rho\ge1$ -- donne $\pi(\Inc)\to1$ ou $\pi(\Inc)=1/2$)
s'applique mot pour mot.
\end{proof}

\begin{proof}[Preuve de (iii)$\Rightarrow$(ii), par contraposée]
Supposons (ii) faux : il existe au moins un littéral en défaut. Par les
Lemmes~\ref{lem:witness} et~\ref{lem:witness-neg} (appliqués au littéral en
question, quel que soit son rôle $J^+$, $J^-$, ou non pertinent), on obtient
\[
  \lim_{N\to\infty}\limsup_{t\to\infty}P(C_t\ne C_{\mathrm{true}}) \ge \frac12 \ne 0,
\]
donc (iii) est faux. Par contraposée, (iii)$\Rightarrow$(ii).
\end{proof}

\section{Conclusion}

Les trois assertions (i), (ii), (iii) sont bien équivalentes : (i)$\iff$(ii)
est une identité algébrique (rien de neuf), (ii)$\Rightarrow$(iii) est le
Théorème 4 déjà prouvé, et (iii)$\Rightarrow$(ii) est nouveau -- démontré
ici par contraposée, littéral par littéral, en distinguant le cas où le
signe est carrément inversé ($\rho\ne1$, le mauvais régime domine par le
Théorème 1) du cas limite où il est exactement nul ($\rho=1$, loi
stationnaire uniforme, probabilité d'erreur bloquée à $1/2$ quel que soit
$N$). Florent avait raison : c'est la même stratégie que pour le Corollaire 2
(exploiter le fait qu'une hypothèse ``pour toute loi'' ou ``ça marche
toujours'' peut être mise en défaut par un seul témoin bien choisi), appliquée
à un théorème différent.

\end{document}
